\appendix
\section{实现细节与复现边界}
\label{app:implementation}
以下细节来自当前可读取的 V5 驱动脚本与配套评价代码，用于复现现有实现，不表示所有历史输出均由同一文件版本生成。最终提交应将脚本提交号、完整依赖、模型权重指纹、运行配置和结果清单绑定到同一版本。

\paragraph{候选清理阈值.}
设前景像素数为 $N_F$、主连通域面积为 $a_0$。附属连通域保留阈值为 $\max(96,\lfloor0.18a_0\rfloor)$；小孔填充面积上限为 $\max(48,\lfloor0.00015N_F\rfloor)$，不填与图像边界连接的背景区域。清理后候选最小面积为 $\max(a_{\min},96,\lfloor0.00012N_F\rfloor)$；覆盖仲裁后的阈值使用 $\max(a_{\min},96,\lfloor0.00010N_F\rfloor)$。其中 $a_{\min}$ 是运行配置中的用户阈值，最终取值为 \TBD{冻结配置}。这些像素阈值依赖输入分辨率，不能移至另一分辨率后仍宣称处理相同。

\paragraph{反馈与记录.}
每个初始成功候选至多进行一次反馈重试。首轮有效结果作为反馈失败时的回退；反馈修改后的框和点用于重试，但本版后处理仍引用原始规划记录。应同时归档原始响应、实际送入 SAM2 的提示、候选掩码、清理掩码、最终标签和每个部件的淘汰原因。部分日志面积字段产生于残余像素填充之前，最终面积统计应直接从导出的整数标签图重新计算，不能混用中间计数。

\paragraph{已核验与待核验项.}
当前已核验的代码覆盖二维规划调用、候选反馈、清理融合、标签记录和二维匹配评价。规划驱动依赖的完整提示模块、最终全量运行版本、模型权重指纹，以及每个对象的 GLB/VXZ 身份仍需在冻结清单中补齐。论文图中的共享资产约束须由清单证实；不同历史运行的日志不能拼接为一次一致实验。当前没有可据以宣称完成三维逐面、逐顶点或逐体素语义 ID 绑定的验证记录。

\section{实验补全清单}
\label{app:pending}
完成主实验前，需确定开发／测试对象 ID、数据授权与人工参考制作流程，确认 D 参考与生成几何在源视角上的可评价关系，并明确无效样本的处理规则。主比较还需固定各基线的粒度、名称输入和自动提示协议，尤其记录 GeoSAM2 的共享引导适配方式。全部数字以对象级结果为依据，表中每个聚合值应可追溯到分子、分母和结果文件。

人工评价需补充评分者人数、盲评方式、评分锚点、一致性统计和名称核验记录。效率实验需补充硬件、环境、计时起止点、预热及重复设置、缓存策略和坍塌判据。可视化应从同版本配对结果中选取；柱状图与折线图待所需独立变量与对象级结果齐备后绘制。本版空白图框用于预留版面，不表示相应实验已经完成。

\section{图表来源与版本说明}
三张现有插图由作者提供，按指定顺序作为方法图、数据处理图和动机图导入；原始图像未被改画。图\ref{fig:data}中 D 分支的英文描述与“人工直接制作并渲染”的口径仍需在最终图源中统一。图\ref{fig:motivation}明确为概念示意，不属于定性实验。三张表的行、条件及指标来自本次提供的工作簿，全部数值保持 TBD；不使用此前 V3/V4 的成功率、SAM2 自预测质量分数或少量样本观察补填。本文献表依据核验到的论文原始来源编制，预印本条目不虚构会议或期刊录用状态。
